\subsection*{Problems}

\textbf{14.1.} Continue with the probabilistic model of Polya's urn in Problem 13.11F o r

i = 1 ,2,3,... ,l e tXi and Yi , respectively, be numbers of white and black balls

drawn in the first i steps. They satisfy Xi \geq 0, Yi \geq 0, and Xi + Yi = i for all i:

(a) Show that, for k = 1,2,3,... ,i, the probability of Xi = k is

P(Xi =k) =

( i

k

) w(w+ 1 )\cdot\cdot\cdot (w+ k - 1 )b(b+ 1 )\cdot\cdot\cdot (b+ i - k - 1 )

(b+ w)(b + w + 1 )\cdot\cdot\cdot (b+ w + i - 1 ) .

(b) Hence, show that Xi/i converges in distribution to Beta(w, b).

(Hint: Use Stirling's approximation and ( x + a)/ (x +b) xa-b.)

\textbf{14.2.} For n = 1 , 2, 3,... ,l e t Xn denote a random variable with distribution

( n \alpha , \beta)Prove that

1\sqrt n\alpha\beta

(

Xn -n\alpha\beta

) D

-\to N( 0,1).

\textbf{14.3.} Let Xn be the random variable distributed according to the Gaussian distribu-

tion N( 0,n ),f o r n = 1, 2,3,... Show that this sequence of random variables fails

to satisfy all four conditions in Levy's continuity theorem.

\textbf{14.4.} For n = 1, 2, 3,... ,l e t Xn be a Gaussian random variable with distribution

N( 0,1/n)Show that Xn converges in distribution, and find the limit distribution.

\textbf{14.5.} Take Z as the sample space with the discrete topology \mathcal{F}= 2Z. Prove that

a sequence of probability measures converges weakly if and only if it converges in

total variation distance.

\textbf{14.6.} Given an arbitrary sequence of random variables (Xn)\infty

n=1. Show that we can

find positive constants cn,f o rn \geq1, such that cnXn converges weakly to 0.

\textbf{14.7.} Let (Xn)\infty

n=1 and (Yn)\infty

n=1 be two sequences of random variables such that

Xn and Yn are independent random variables defined on the same probability space.

Show that if Xn

D

\to X and Yn

D

\to Y , then Xn +Yn converges to X+Y in distribution

(cf. Problem 10.10). (Hint: Apply Levy's continuity theorem.)

\textbf{14.8.} Let (X n)\infty

n=1 be a sequence of random variables satisfying the following

properties:

(i) For all n , the moment generating function mXn (t) of Xn is defined for t \in

[-\delta, \delta].

(ii) For all t \in[ - \delta, \delta], the moment generating functions mXn (t) converge to the

moment generating function mX(t) of random variable X:

(a) Prove that the sequence (Xn)n\geq1 is tight.

(b) Prove that Xn converges to X in distribution.

\textbf{14.9.} (Continuous mapping theorem for convergence in distribution) Let Xn be a

sequence of random variables converging to X in distribution. Let f : R \to R be a

function that is continuous at the set Sf \subseteq R. Show that if P(X \in Sf ) = 1, then

f( Xn)

L

-\to f( X) . (Hint: Apply Levy's continuity theorem.)

\textbf{14.10.} Consider a sequence of random variables (X n)\infty

n=1 that is uniformly

bounded, ie., there exists a constant c such that \vertXn\vert\leq c as. for all n Let X

be a random variable that is also bounded with probability 1. Prove that Xn

D

\to X

if and only if lim n\to\infty E[Xk

n]= E[X] for all k \geq 1. (Hint: Use the result in

Problem 14.8.)

\textbf{14.11.} Generalize Example 14.4.3 to the case when m/n approaches a fixed

constant r,f o r0 <r < 1.

\textbf{14.12.} (Uniform integrability) A sequence of random variables (Xn)\infty

n=1 is said to

be uniformly integrable if for every \epsilon> 0, there is a sufficiently large M such that

E[Xn1\vertXn\vert\geqM]\leq \epsilonfor all n. Prove the followings:

(a) A uniformly integrable sequence of random variables is tight.

(b) If a sequence of random variables is uniformly integrable and is converging in

probability, then it is converging in the mean.
